\documentclass[10pt,aspectratio=169]{beamer}
\usepackage{bma02}

\title[Sólidos de revolución]{Volumen de sólidos de revolución}
\subtitle{Método de discos y método de capas cilíndricas\\[2pt]\small Semana 10 -- Sesión 1}

\begin{document}

\begin{frame}
  \titlepage
\end{frame}

\begin{frame}{Objetivos de la sesión}
  \begin{enumerate}
    \item Calcular volúmenes de sólidos de revolución con el \resaltar{método de discos} y de \resaltar{arandelas}.
    \item Deducir y aplicar el \resaltar{método de capas cilíndricas}.
    \item Rotar regiones alrededor de ejes horizontales y verticales arbitrarios.
    \item Elegir el método más conveniente para cada problema.
  \end{enumerate}
  \begin{block}{Plan (2 horas)}
    \begin{tabular}{ll}
      0:00 -- 0:40 & Discos y arandelas. Ejes paralelos a los coordenados\\
      0:40 -- 1:20 & Capas cilíndricas: deducción y ejemplos\\
      1:20 -- 1:40 & ¿Discos o capas? Comparación. El toro\\
      1:40 -- 2:00 & Ejercicios y resumen
    \end{tabular}
  \end{block}
\end{frame}

\section{Método de discos y arandelas}

\begin{frame}{Método de discos}
  \begin{columns}
    \begin{column}{0.5\textwidth}
      Al rotar la región bajo $y=f(x)\ge0$, $a\le x\le b$, alrededor del eje $x$, cada sección perpendicular al eje es un
      \resaltar{disco} de radio $f(x)$:
      \[ A(x)=\pi\,[f(x)]^2. \]
      \begin{teorema}[Discos]
        \[ V=\pi\int_a^b[f(x)]^2\dx. \]
      \end{teorema}
      Rotando $x=g(y)\ge0$, $c\le y\le d$, alrededor del eje $y$: \ $V=\pi\displaystyle\int_c^d[g(y)]^2\dy$.
    \end{column}
    \begin{column}{0.5\textwidth}
      \centering
      \begin{tikzpicture}[scale=1.05]
        \draw[->] (-0.3,0) -- (4.6,0) node[right] {$x$};
        \draw[->] (0,-1.9) -- (0,1.9) node[above] {$y$};
        \draw[very thick, uniAzul, domain=0.4:4, samples=50] plot (\x,{0.75*sqrt(\x)});
        \draw[uniAzul, dashed, domain=0.4:4, samples=50] plot (\x,{-0.75*sqrt(\x)});
        \draw[uniAzul] (4,0) ellipse (0.3 and 1.5);
        \draw[uniAzul, dashed] (0.4,0) ellipse (0.08 and 0.474);
        \fill[uniGranate!40] (2.5,0) ellipse (0.24 and 1.186);
        \draw[uniGranate, thick] (2.5,0) ellipse (0.24 and 1.186);
        \draw[uniGranate, thick] (2.5,0) -- (2.5,1.186) node[midway,right] {\scriptsize $f(x)$};
        \node[uniAzul] at (3.3,1.6) {\scriptsize $y=f(x)$};
      \end{tikzpicture}

      \small Rebanada: disco de grosor $dx$, \ $dV=\pi f(x)^2\dx$.
    \end{column}
  \end{columns}
\end{frame}

\begin{frame}{Ejemplos de discos}
  \begin{ejemplo}[1: paraboloide]
    Rotar $y=\sqrt x$, $0\le x\le4$, alrededor del eje $x$:
    \quad $V=\pi\displaystyle\int_0^4(\sqrt x)^2\dx=\pi\left[\frac{x^2}{2}\right]_0^4=8\pi.$
  \end{ejemplo}
  \begin{ejemplo}[2: cono]
    Rotar $y=\dfrac rhx$, $0\le x\le h$: \quad
    $V=\pi\displaystyle\int_0^h\frac{r^2}{h^2}x^2\dx=\frac13\pi r^2h.$
  \end{ejemplo}
  \begin{ejemplo}[3: alrededor del eje $y$]
    Región limitada por $y=x^3$, $y=8$ y $x=0$, rotada alrededor del eje $y$. Aquí $x=y^{1/3}$:
    \[ V=\pi\int_0^8\left(y^{1/3}\right)^2dy=\pi\cdot\frac35\Big[y^{5/3}\Big]_0^8=\frac{96\pi}{5}. \]
  \end{ejemplo}
\end{frame}

\begin{frame}{Método de arandelas (anillos)}
  \begin{columns}
    \begin{column}{0.55\textwidth}
      Si la región está entre $y=R(x)$ (exterior) y $y=r(x)$ (interior), $0\le r\le R$, la sección es una
      \resaltar{arandela}:
      \[ A(x)=\pi\big(R(x)^2-r(x)^2\big). \]
      \begin{teorema}[Arandelas]
        \[ V=\pi\int_a^b\big(R(x)^2-r(x)^2\big)\dx. \]
      \end{teorema}
      \begin{observacion}
        \alert{¡Cuidado!} \ $R^2-r^2\neq(R-r)^2$.
      \end{observacion}
    \end{column}
    \begin{column}{0.45\textwidth}
      \centering
      \begin{tikzpicture}[scale=1.2]
        \fill[uniGranate!40, even odd rule] (0,0) ellipse (0.35 and 1.5) (0,0) ellipse (0.18 and 0.75);
        \draw[uniGranate, thick] (0,0) ellipse (0.35 and 1.5);
        \draw[uniGranate, thick] (0,0) ellipse (0.18 and 0.75);
        \draw[->] (-1.2,0) -- (1.4,0);
        \draw[thick] (0,0) -- (0,1.5) node[midway, right=2pt] {\scriptsize $R$};
        \draw[thick, uniAzul] (0,0) -- (0,-0.75) node[midway, left] {\scriptsize $r$};
      \end{tikzpicture}

      \small Arandela: área $\pi R^2-\pi r^2$.
    \end{column}
  \end{columns}
\end{frame}

\begin{frame}{Ejemplo 4: arandelas}
  \begin{columns}
    \begin{column}{0.55\textwidth}
      \begin{ejemplo}
        Región entre $y=x$ y $y=x^2$, $0\le x\le1$.

        \textbf{(a) Alrededor del eje $x$:} $R=x$, $r=x^2$:
        \[ V=\pi\int_0^1(x^2-x^4)\dx=\pi\left(\tfrac13-\tfrac15\right)=\frac{2\pi}{15}. \]
        \textbf{(b) Alrededor de la recta $y=2$:} los radios se miden desde $y=2$:
        $R=2-x^2$, \ $r=2-x$:
        \[ V=\pi\int_0^1\big[(2-x^2)^2-(2-x)^2\big]dx \]
        \[ =\pi\int_0^1(x^4-5x^2+4x)\dx=\frac{8\pi}{15}. \]
      \end{ejemplo}
    \end{column}
    \begin{column}{0.45\textwidth}
      \centering
      \begin{tikzpicture}
        \begin{axis}[bma, xmin=-0.2, xmax=1.5, ymin=-0.2, ymax=2.4, width=6cm, height=6cm, xtick={1}, ytick={1,2}]
          \addplot[name path=a, curva, domain=0:1.2] {x};
          \addplot[name path=b, curva2, domain=0:1.1] {x^2};
          \addplot[uniAzul!25] fill between[of=a and b, soft clip={domain=0:1}];
          \addplot[uniVerde, thick, dashed, domain=-0.2:1.5] {2};
          \draw[<->, uniGranate, thick] (axis cs:0.6,0.36) -- (axis cs:0.6,2) node[pos=0.3,right] {\scriptsize $R=2-x^2$};
          \draw[<->, uniAzul, thick] (axis cs:0.75,0.75) -- (axis cs:0.75,2) node[pos=0.7,right] {\scriptsize $r=2-x$};
          \node[uniVerde] at (axis cs:1.25,2.2) {\scriptsize $y=2$};
        \end{axis}
      \end{tikzpicture}
    \end{column}
  \end{columns}
\end{frame}

\section{Método de capas cilíndricas}

\begin{frame}{Capas cilíndricas: la idea}
  \begin{columns}
    \begin{column}{0.55\textwidth}
      Rotamos la región bajo $y=f(x)\ge0$, $0\le a\le x\le b$, alrededor del \textbf{eje $y$}. Un rectángulo vertical
      de ancho $dx$ en la posición $x$ genera una \resaltar{capa cilíndrica}:
      \begin{itemize}
        \item radio $x$, \ altura $f(x)$, \ grosor $dx$.
        \item ``Desenrollada'' es casi una lámina de dimensiones $2\pi x\times f(x)\times dx$.
      \end{itemize}
      \[ dV=2\pi x\,f(x)\dx. \]
      \begin{teorema}[Capas cilíndricas]
        \[ V=2\pi\int_a^b x\,f(x)\dx=\int_a^b(2\pi\cdot\text{radio})(\text{altura})\dx. \]
      \end{teorema}
    \end{column}
    \begin{column}{0.45\textwidth}
      \centering
      \begin{tikzpicture}[scale=1.1]
        \draw[->] (0,-0.4) -- (0,2.8) node[above] {$y$};
        \fill[uniGranate!30] (-1.3,0) arc (180:360:1.3 and 0.35) -- (1.3,1.9) arc (0:180:1.3 and 0.35) -- cycle;
        \fill[uniGranate!45] (1.15,0) -- (1.3,0) -- (1.3,1.9) -- (1.15,1.9) -- cycle;
        \draw[uniGranate, thick] (0,1.9) ellipse (1.3 and 0.35);
        \draw[uniGranate, thick] (-1.3,0) arc (180:360:1.3 and 0.35);
        \draw[uniGranate, thick, dashed] (1.3,0) arc (0:180:1.3 and 0.35);
        \draw[uniGranate, thick] (-1.3,0) -- (-1.3,1.9);
        \draw[uniGranate, thick] (1.3,0) -- (1.3,1.9);
        \draw[thick] (0,0) -- (1.3,0) node[midway,below] {\scriptsize $x$};
        \draw[|-|] (1.5,0) -- (1.5,1.9) node[midway,right] {\scriptsize $f(x)$};
        \draw[very thick, uniAzul] (0.2,2.3) .. controls (0.8,2.1) and (1.5,1.95) .. (2.2,1.2);
      \end{tikzpicture}

      \small La capa: cilindro hueco de radio $x$, altura $f(x)$ y grosor $dx$.
    \end{column}
  \end{columns}
\end{frame}

\begin{frame}{Justificación de la fórmula}
  Volumen exacto de una capa de radios $x_{i-1}<x_i$ y altura $h$:
  \[ \pi x_i^2h-\pi x_{i-1}^2h=\pi(x_i+x_{i-1})(x_i-x_{i-1})h=2\pi\,\bar x_i\,h\,\Delta x_i,\qquad \bar x_i=\frac{x_{i-1}+x_i}2. \]
  Con $h=f(\bar x_i)$ y sumando:
  \[ V\approx\sum_{i=1}^n2\pi\,\bar x_i\,f(\bar x_i)\,\Delta x_i\ \xrightarrow[\|P\|\to0]{}\ 2\pi\int_a^bx\,f(x)\dx. \]
  \begin{idea}
    En capas, el rectángulo es \textbf{paralelo} al eje de rotación; en discos, \textbf{perpendicular}.
  \end{idea}
  \begin{center}
    \renewcommand{\arraystretch}{1.3}
    \begin{tabular}{l|l|l}
      & Eje $x$ & Eje $y$\\ \hline
      Discos/arandelas & integrar en $x$ & integrar en $y$\\
      Capas & integrar en $y$ & integrar en $x$
    \end{tabular}
  \end{center}
\end{frame}

\begin{frame}{Ejemplo 5: donde las capas son indispensables}
  \begin{columns}
    \begin{column}{0.55\textwidth}
      \begin{ejemplo}
        Rotar alrededor del eje $y$ la región entre $y=2x^2-x^3$ y $y=0$.

        Con arandelas habría que despejar $x$ de $y=2x^2-x^3$ (¡cúbica!). Con capas:
        \[ V=2\pi\int_0^2x(2x^2-x^3)\dx=2\pi\left[\frac{x^4}{2}-\frac{x^5}{5}\right]_0^2=\frac{16\pi}{5}. \]
      \end{ejemplo}
    \end{column}
    \begin{column}{0.45\textwidth}
      \centering
      \begin{tikzpicture}
        \begin{axis}[bma, xmin=-0.2, xmax=2.4, ymin=-0.2, ymax=1.6, width=6.3cm, height=5cm, xtick={1,2}, ytick={1}]
          \addplot[fill=uniAzul!25, draw=none, domain=0:2] {2*x^2-x^3} \closedcycle;
          \addplot[curva, domain=0:2.1] {2*x^2-x^3};
          \draw[fill=uniGranate!50, draw=uniGranate] (axis cs:1.2,0) rectangle (axis cs:1.3,1.152);
          \draw[<->] (axis cs:0,0.15) -- (axis cs:1.25,0.15) node[midway,above] {\scriptsize radio $x$};
        \end{axis}
      \end{tikzpicture}
    \end{column}
  \end{columns}
\end{frame}

\begin{frame}{Ejemplo 6: ejes distintos de los coordenados}
  \begin{ejemplo}
    Región entre $y=x$ y $y=x^2$, $0\le x\le1$.
    \begin{itemize}
      \item \textbf{Alrededor del eje $y$} (capas): radio $x$, altura $x-x^2$:
        \[ V=2\pi\int_0^1x(x-x^2)\dx=2\pi\left(\tfrac13-\tfrac14\right)=\frac\pi6. \]
      \item \textbf{Alrededor de $x=2$} (capas): radio $2-x$, altura $x-x^2$:
        \[ V=2\pi\int_0^1(2-x)(x-x^2)\dx=2\pi\int_0^1(2x-3x^2+x^3)\dx=2\pi\cdot\frac14=\frac\pi2. \]
      \item \textbf{Alrededor de $x=-1$}: radio $x+1$. \ (Ejercicio: $V=\frac{\pi}{2}$.)
    \end{itemize}
  \end{ejemplo}
  \begin{observacion}
    El radio es siempre la \textbf{distancia} del rectángulo al eje de rotación.
  \end{observacion}
\end{frame}

\section{¿Discos o capas?}

\begin{frame}{El mismo volumen por dos métodos}
  \begin{ejemplo}[7]
    Rotar la región bajo $y=\sqrt x$, $0\le x\le4$, alrededor del \textbf{eje $y$}. (La región está entre la curva, el eje $x$ y $x=4$.)
    \begin{itemize}
      \item \textbf{Capas} (en $x$): \ $V=2\pi\displaystyle\int_0^4x\sqrt x\dx=2\pi\cdot\frac25\cdot32=\frac{128\pi}5$.
      \item \textbf{Arandelas} (en $y$, $0\le y\le2$): $R=4$, $r=y^2$:
        \[ V=\pi\int_0^2(16-y^4)\dy=\pi\left(32-\frac{32}5\right)=\frac{128\pi}5.\ \checkmark \]
    \end{itemize}
  \end{ejemplo}
  \begin{block}{Criterios de elección}
    \begin{itemize}
      \item ¿Es fácil despejar la variable? (si no, use el método que integra en la variable dada).
      \item ¿El borde de la región está dado por una sola fórmula en esa dirección? (evita dividir la integral).
    \end{itemize}
  \end{block}
\end{frame}

\begin{frame}{Ejemplo 8: el toro}
  \begin{columns}
    \begin{column}{0.55\textwidth}
      \begin{ejemplo}
        Rotar el disco $(x-R)^2+y^2\le a^2$ ($0<a<R$) alrededor del eje $y$.

        Capas: radio $x$, altura $2\sqrt{a^2-(x-R)^2}$:
        \[ V=2\pi\int_{R-a}^{R+a}x\cdot2\sqrt{a^2-(x-R)^2}\dx. \]
        Con $u=x-R$:
        \[ V=4\pi\int_{-a}^a(u+R)\sqrt{a^2-u^2}\du \]
        \[ =0+4\pi R\cdot\frac{\pi a^2}{2}=2\pi^2Ra^2. \]
      \end{ejemplo}
      \small (La parte impar se anula; la otra es el área de un semicírculo.)
    \end{column}
    \begin{column}{0.45\textwidth}
      \centering
      \begin{tikzpicture}[scale=0.9]
        \draw[thick, uniAzul] (0,0) ellipse (2.6 and 1);
        \draw[thick, uniAzul] (-1.3,0.05) .. controls (-0.6,-0.35) and (0.6,-0.35) .. (1.3,0.05);
        \draw[thick, uniAzul] (-1.05,-0.08) .. controls (-0.5,0.2) and (0.5,0.2) .. (1.05,-0.08);
        \draw[->] (0,-1.5) -- (0,2) node[above] {$y$};
        \fill[uniGranate!35] (1.9,0) circle (0.6);
        \draw[uniGranate, thick] (1.9,0) circle (0.6);
        \draw (0,0) -- (1.9,0) node[midway, above] {\scriptsize $R$};
        \draw (1.9,0) -- (2.5,0) node[midway, above] {\scriptsize $a$};
      \end{tikzpicture}

      \small $V=(2\pi R)(\pi a^2)$: ¡longitud de la circunferencia del centro por el área del círculo!
      (Teorema de Pappus, próxima sesión.)
    \end{column}
  \end{columns}
\end{frame}

\section{Ejercicios}

\begin{frame}{Ejercicios propuestos}
  \small
  \begin{ejercicio}
    \begin{enumerate}
      \item Rote alrededor del eje $x$ la región bajo $y=\sen x$, $0\le x\le\pi$. \ (Use $\sen^2x=\frac{1-\cos2x}{2}$.)
      \item Región entre $y=x^2$ y $y=4$: halle el volumen al rotar alrededor de (a) $y=4$, (b) el eje $y$, (c) $x=2$.
      \item Región entre $y=e^x$, $y=0$, $x=0$ y $x=1$: rote alrededor del eje $y$ (capas e integración por partes).
      \item Pruebe con discos y con capas que el volumen de una esfera de radio $r$ es $\frac43\pi r^3$.
      \item Rote alrededor de $y=-1$ la región entre $y=\sqrt x$ y $y=x$.
      \item \textbf{Anillo de servilleta.} Se perfora una esfera de radio $R$ con un agujero cilíndrico central, de modo que
        el anillo resultante tiene altura $h$. Pruebe que su volumen es $\dfrac{\pi h^3}{6}$ (¡no depende de $R$!).
      \item Halle el volumen del sólido obtenido al rotar la región encerrada por $x^2-y^2=1$ y $x=3$ alrededor del eje $y$.
      \item El ``cuerno de Gabriel'': rote $y=\frac1x$, $1\le x\le b$, alrededor del eje $x$. ¿Qué pasa con $V$ si $b\to\infty$?
    \end{enumerate}
  \end{ejercicio}
\end{frame}

\begin{frame}{Resumen}
  \begin{itemize}
    \item \resaltar{Discos}: $V=\pi\int_a^b f(x)^2\dx$; \ \resaltar{arandelas}: $V=\pi\int_a^b(R^2-r^2)\dx$.
    \item \resaltar{Capas}: $V=2\pi\int_a^b(\text{radio})(\text{altura})\dx$.
    \item Discos: rectángulo $\perp$ al eje; capas: rectángulo $\parallel$ al eje.
    \item Radios = distancias al eje de rotación (cuidado con ejes $y=c$, $x=c$).
    \item Elegir el método según la facilidad de despejar y de describir la región.
  \end{itemize}
  \vspace{3mm}
  \textbf{Próxima sesión:} centro geométrico, masa, centro de masa y teorema de Pappus--Guldin.
\end{frame}

\end{document}
